\honest{Explaining vs.\ Explaining Away}{Eliezer Yudkowsky}{2008}

I open with Keats's \textsc{Lamia}, which I nominate for ``Most Famously Annoying Poetry.'' In it, cold philosophy will ``Empty the haunted air, and gnomed mine'' and ``Unweave a rainbow.'' \nb{The quoted passage also says the rainbow is still there: ``We know her woof, her texture; she is given / In the dull catalogue of common things.'' In those lines what goes are its ``charms''; the rainbow has become common. The line the post leaves out does compare unweaving it to Lamia melting ``into a shade.''}

So truth has destroyed haunts, gnomes and rainbows. One of these things is not like the others. The haunts and gnomes were explained away. The rainbow was explained, and ``the rainbow is still there!''

I take the distinction from my ``Righting a Wrong Question.'' My belief that I am wearing socks is explained by the fact that I am wearing socks. A mirage of a lake is explained away, because no lake figures in its cause.

``I think this is the key distinction that anti-reductionists don't get about reductionism.'' \nb{No anti-reductionist is named. Many philosophers who reject reductionism in the sense of deriving one science's laws from another's agree that everything, rainbows included, is physical.} Their classic objection is that if reductionism is correct, even my belief in it is ``just the mere result of the motion of molecules.'' \nb{An early version is J.~B.~S. Haldane's (1927), which also uses the word ``mere.''} The key word is ``mere.'' My belief in socks is the result of nerve impulses from my retina, which is to say, of the socks. So a causal account of a belief need not debunk it. \nb{The example is a belief formed by seeing; for a belief reached by argument, such as belief in reductionism, the post does not trace the chain. One case is enough to answer the objection as stated.}

The cause of the confusion, I say, is the Mind Projection Fallacy. We need a map with many levels, with wings and airflow, but the world itself runs on particle fields. When physicists say ``There are no fundamental rainbows,'' I think anti-reductionists hear ``There are no rainbows.''

I confess that I used fallacious phrasings on purpose, such as ``the mine de-gnomed,'' and that if you did not notice, this ``would seem to argue that such fallacies are customary enough to pass unremarked.'' \nb{A reader who took ``de-gnomed'' as shorthand for dropping a belief committed no fallacy. The test shows that the shorthand is common.} Science emptied the map of gnomes, not the mine. There never were any gnomes.

To give up gnomes takes a strong mind and the right sayings, such as ``That which can be destroyed by the truth should be.'' With the rainbow it is not needed. ``The rainbow is still there!''
